\documentclass[12pt,a4paper]{article}
\usepackage[utf8]{inputenc}
\usepackage[T2A]{fontenc}
\usepackage[russian]{babel}
\usepackage{textcomp}
\usepackage{amsmath, amssymb, amsthm}
\usepackage{graphicx}
\usepackage{hyperref}
\usepackage{listingsutf8}
\usepackage{xcolor}
\usepackage{geometry}
\usepackage{float}
\usepackage{caption}
\usepackage{subcaption}
\usepackage{enumitem}
\usepackage{titlesec}
\usepackage{fancyhdr}
\usepackage{array}
\usepackage{longtable}
\usepackage{booktabs}

% Настройка полей
\geometry{
    a4paper,
    left=2.5cm,
    right=1.5cm,
    top=2cm,
    bottom=2cm
}

% Настройка заголовков
\titleformat{\section}{\Large\bfseries}{\thesection}{1em}{}
\titleformat{\subsection}{\large\bfseries}{\thesubsection}{1em}{}
\titleformat{\subsubsection}{\normalsize\bfseries}{\thesubsubsection}{1em}{}

% Настройка кода (в листингах только латиница)
\lstset{
    language=Python,
    basicstyle=\ttfamily\small,
    keywordstyle=\color{blue},
    commentstyle=\color{green!40!black},
    stringstyle=\color{red},
    numbers=left,
    numberstyle=\tiny,
    breaklines=true,
    frame=single,
    captionpos=b,
    showstringspaces=false,
    tabsize=4
}

% Настройка колонтитулов
\pagestyle{fancy}
\fancyhf{}
\fancyhead[L]{Лабораторная работа: Проектирование AI-сервиса на FastAPI}
\fancyhead[R]{\thepage}
\fancyfoot[C]{}

\begin{document}

\begin{titlepage}
    \begin{center}
        \textbf{МИНИСТЕРСТВО НАУКИ И ВЫСШЕГО ОБРАЗОВАНИЯ РОССИЙСКОЙ ФЕДЕРАЦИИ}
        \vspace{0.5cm}

        \textbf{Федерального государственного автономного образовательного учреждения \\ высшего образования}
        \vspace{0.3cm}

        \textbf{«НАЦИОНАЛЬНЫЙ ИССЛЕДОВАТЕЛЬСКИЙ ТЕХНОЛОГИЧЕСКИЙ УНИВЕРСИТЕТ МИСИС»}
        \vspace{2cm}

        \textbf{Институт компьютерных наук}
        \vspace{0.3cm}

        \textbf{Кафедра Инженерной кибернетики}
        \vspace{1cm}

        \textbf{Лабораторная работа №1}
        \vspace{0.3cm}

        \textbf{по дисциплине «Современные интеллектуальные сетевые сервисы»}
        \vspace{1.5cm}

        \huge{\textbf{Проектирование AI-сервиса \\ на FastAPI}}
        \vspace{2cm}
    \end{center}

    \begin{flushright}
        \textbf{Выполнил:} \\
        Студент группы МПИ 25-1-2\\
        Д.А. Комарова \\
        \vspace{0.5cm}
        \textbf{Проверил:} \\
        Преподаватель \\
        М.Е. Корнет \\
        \vspace{2cm}
        \textbf{Оценка:} \hrulefill
    \end{flushright}

    \begin{center}
        \vfill
        Москва, 2026
    \end{center}
\end{titlepage}

\tableofcontents
\newpage

\section{Цель работы}

Цель практической работы — научиться проектировать программную архитектуру AI-сервиса, использующего модель машинного обучения, определить место API в общей ИИ-системе и реализовать REST API на FastAPI, изучив основные принципы программного взаимодействия между AI-моделью и внешними информационными системами.

\section{Задачи работы}

В ходе работы необходимо:

\begin{itemize}
    \item определить потребителей AI-сервиса;
    \item определить формат входных и выходных данных;
    \item спроектировать контур инференса;
    \item разделить синхронные и асинхронные операции;
    \item определить основные программные компоненты системы;
    \item выбрать протоколы взаимодействия;
    \item сформировать структурированный контракт REST API для AI-сервиса;
    \item определить способы хранения результатов;
    \item определить требования к логированию и наблюдаемости;
    \item выявить основные архитектурные риски;
    \item создать приложение на FastAPI и реализовать REST endpoints;
    \item изучить методы GET и POST, передачу данных через JSON, использование path- и query-параметров;
    \item описывать входные и выходные схемы с помощью Pydantic и выполнять валидацию входных данных;
    \item возвращать корректные HTTP-коды и обрабатывать ошибки;
    \item изучить автоматически создаваемую OpenAPI-документацию и протестировать API через Swagger UI;
    \item реализовать минимальный прототип API на FastAPI.
\end{itemize}

\section{Исходные данные}

Банк разрабатывает сервис агроскоринга. Система должна оценивать риск при работе с сельскохозяйственными предприятиями. Пользователь или другая информационная система передаёт сведения о хозяйстве: регион, площадь посевов, тип основной культуры, погодные показатели, историю платежей, количество просроченных платежей, размер текущей задолженности.

На основании этих данных модель машинного обучения должна определить риск профиля предприятия. Результат должен содержать числовую оценку риска, категорию риска и рекомендацию оператору банка.

Пример ответа:

\begin{lstlisting}[language=Python, caption={Пример ответа /predict}]
{
    "risk_score": 0.73,
    "risk_level": "high",
    "recommendation": "Additional check required"
}
\end{lstlisting}

Сервис должен позволять:

\begin{itemize}
    \item проверить, работает ли API;
    \item получить информацию о модели;
    \item выполнить оценку риска одного хозяйства;
    \item получить сохранённый результат прогноза;
    \item получить несколько последних прогнозов;
    \item корректно сообщать клиенту об ошибках.
\end{itemize}

Предполагается, что в будущем этим API будут пользоваться: CRM банка, веб-портал, мобильное приложение, внутренние банковские сервисы, BI-система.

\section{Требуемые endpoints}

API должен иметь следующие endpoints:

\begin{table}[H]
\centering
\caption{Список endpoints API}
\label{tab:endpoints}
\begin{tabular}{|l|l|l|}
\hline
\textbf{Метод} & \textbf{Endpoint} & \textbf{Назначение} \\
\hline
GET  & /health                    & Проверка работоспособности API \\
GET  & /model-info                & Информация о модели \\
POST & /predict                   & Оценка риска хозяйства \\
GET  & /predictions/\{request\_id\} & Получение прогноза по ID \\
GET  & /predictions               & Список прогнозов \\
\hline
\end{tabular}
\end{table}

\section{Теоретическая часть}

\subsection{Что такое REST API}

REST API — способ взаимодействия программных систем через HTTP. Клиент отправляет запрос (HTTP Request), сервер обрабатывает его и возвращает ответ (HTTP Response).

Клиенту не нужно знать: на каком языке написана модель, где она находится, является ли модель scikit-learn или PyTorch, сколько функций выполняется внутри сервера. Клиенту достаточно знать контракт API.

\subsection{Место API в интеллектуальной системе}

AI-сервис редко существует изолированно. Как правило, он встроен в более крупную информационную систему. Типовая схема:

\begin{center}
Пользователь / UI $\rightarrow$ Внешняя система (CRM, ERP, мобильное приложение) \\
$\rightarrow$ HTTP / JSON $\rightarrow$ API-слой (FastAPI) \\
$\rightarrow$ Контур инференса: предобработка $\rightarrow$ модель $\rightarrow$ постобработка \\
$\rightarrow$ Хранилище результатов: БД, лог, очередь
\end{center}

API в этой схеме выполняет три роли:

\begin{enumerate}
    \item \textbf{Точка входа} — единый интерфейс для всех внешних потребителей.
    \item \textbf{Изолятор} — скрывает детали реализации модели (фреймворк, версию, способ загрузки).
    \item \textbf{Контракт} — фиксирует форматы входных и выходных данных, коды ответов и правила обработки ошибок.
\end{enumerate}

Благодаря этому модель можно заменить, переобучить или перенести на другой сервер, не меняя код клиентов.

\subsection{Синхронные и асинхронные операции}

В AI-сервисе операции делятся на два класса.

\textbf{Синхронные} — клиент ждёт ответ в рамках одного запроса. Подходят для быстрых моделей (например, инференс за 50--500 мс).

\textbf{Асинхронные} — клиент получает идентификатор задачи и забирает результат позже. Нужны для тяжёлых моделей, пакетной обработки или длительных вычислений.

Выбор между ними — архитектурное решение, влияющее на таймауты, нагрузку и требования к хранилищу.

\subsection{Формат данных: JSON}

Обмен данными между клиентом и AI-сервисом выполняется в формате JSON. Это текстовый формат, который легко читается человеком, поддерживается всеми языками программирования и удобен для различных структур данных (объекты, массивы, числа, строки, булевы значения).

\subsection{HTTP-методы}

В этой работе используются два основных метода. \textbf{GET} используется преимущественно для получения информации. \textbf{POST} используется для передачи данных серверу и создания результата.

\subsection{Структура HTTP-запроса}

Упрощённо HTTP-запрос можно представить следующим образом:

\begin{lstlisting}[language=bash, caption={Структура HTTP-запроса}]
METHOD + URL
HEADERS
BODY
\end{lstlisting}

Например:

\begin{lstlisting}[language=bash, caption={Пример POST-запроса}]
POST /predict
Content-Type: application/json

{
    "farm_id": "FARM-100",
    "area_ha": 2500
}
\end{lstlisting}

\subsection{Path- и query-параметры}

Данные можно передавать не только в теле запроса, но и прямо в URL.

\textbf{Path-параметры} — часть пути, идентифицируют конкретный ресурс, например: \texttt{GET /predictions/123}.

\textbf{Query-параметры} — идут после знака \texttt{?}, обычно используются для фильтрации, сортировки, пагинации: \texttt{GET /predictions?limit=10\&offset=0\&region=Krasnodar}.

Различие важно для проектирования контракта: path-параметр описывает \emph{что} запрашивается, query-параметр уточняет \emph{как} это подаётся (сколько, в каком порядке, с какими фильтрами).

\subsection{Валидация данных и Pydantic}

Клиент может прислать некорректные данные: пропустить обязательное поле, передать строку вместо числа, указать отрицательную площадь.

Чтобы сервис не падал и не выполнял инференс на ошибочных данных, входные и выходные структуры описываются явно — с помощью Pydantic-моделей:

\begin{lstlisting}[caption={Pydantic-модели запроса и ответа}]
from pydantic import BaseModel, Field

class PredictRequest(BaseModel):
    farm_id: str
    region: str
    crop_type: str
    area_ha: float = Field(gt=0)

class PredictResponse(BaseModel):
    risk_score: float
    risk_level: str
    model_version: str
\end{lstlisting}

Pydantic автоматически проверяет типы и ограничения и возвращает понятную ошибку при несоответствии. Это одновременно и защита сервиса, и документация контракта.

\subsection{HTTP-коды ответов}

Сервер сообщает о результате не только телом ответа, но и HTTP-кодом. Основные группы:

\begin{table}[H]
\centering
\caption{Основные HTTP-коды}
\label{tab:http_codes}
\begin{tabular}{|l|l|l|}
\hline
\textbf{Код} & \textbf{Значение} & \textbf{Пример в AI-сервисе} \\
\hline
200 & OK                    & успешный инференс \\
201 & Created               & создана асинхронная задача \\
400 & Bad Request           & некорректное тело запроса \\
404 & Not Found             & задача с таким id не найдена \\
422 & Unprocessable Entity  & ошибка валидации Pydantic \\
500 & Internal Server Error & сбой модели или кода \\
503 & Service Unavailable   & модель ещё не загружена \\
\hline
\end{tabular}
\end{table}

Корректные коды — часть контракта: клиент по ним понимает, нужно ли повторять запрос, исправлять данные или сообщать об ошибке пользователю.

\subsection{Обработка ошибок}

Ошибки бывают трёх типов. Ошибки клиента — некорректный JSON, отсутствует поле, неверный тип; возвращается 400 или 422 с описанием проблемы. Ошибки бизнес-логики — например, запрошен прогноз для несуществующего хозяйства; возвращается 404 или 409. Внутренние ошибки — падение модели, недоступность БД; возвращается 500, а детали пишутся в лог, но не отдаются клиенту.

\subsection{OpenAPI и Swagger UI}

FastAPI автоматически строит OpenAPI-спецификацию — машиночитаемое описание всех endpoints, схем запроса и ответа, кодов ошибок. На её основе генерируется Swagger UI — интерактивная страница, доступная по адресу \texttt{/docs}, где можно увидеть все endpoints, посмотреть схемы Pydantic, отправить тестовый запрос прямо из браузера и получить ответ и код. Это одновременно и документация для разработчиков, и инструмент тестирования, и основа для генерации клиентов на других языках.

\subsection{Логирование и наблюдаемость}

Для того чтобы обеспечить воспроизводимость, отследить деградацию модели и проверять, какая версия модели дала конкретный ответ, фиксируют входящие запросы (метод, путь, идентификатор клиента, время), результат инференса (версия модели, ключевые признаки, предсказание), ошибки (тип, стек, входные данные без персональных данных) и метрики (задержка, доля ошибок, количество запросов).

\subsection{Архитектурные риски}

Даже простой AI-сервис сталкивается с типовыми рисками:

\begin{table}[H]
\centering
\caption{Архитектурные риски и меры}
\label{tab:risks}
\begin{tabular}{|p{4cm}|p{5cm}|p{5cm}|}
\hline
\textbf{Риск} & \textbf{Проявление} & \textbf{Мера} \\
\hline
Долгий инференс & таймауты у клиента & асинхронная обработка, очередь \\
Некорректные данные & падение модели & валидация Pydantic \\
Смена версии модели & непредсказуемые ответы & версионирование API и модели \\
Перегрузка & рост задержек & лимиты, кэширование, горизонтальное масштабирование \\
Отсутствие логов & невозможность разбора инцидента & структурированное логирование \\
Утечка данных & нарушение приватности & минимизация логируемых полей \\
\hline
\end{tabular}
\end{table}

Их выявление на этапе проектирования дешевле, чем устранение в проде.

\section{Создание проекта}

Создайте каталог \texttt{agro\_api/}. Базовая структура:

\begin{lstlisting}[language=bash, caption={Базовая структура проекта}]
agro_api/
    main.py
    schemas.py
    model.py
    requirements.txt
\end{lstlisting}

\section{Создание виртуального окружения}

Рекомендуется создать отдельное окружение Python:

\begin{lstlisting}[language=bash]
python -m venv venv
\end{lstlisting}

Активация:

\begin{lstlisting}[language=bash]
venv\Scripts\activate
\end{lstlisting}

Установите зависимости:

\begin{lstlisting}[language=bash]
pip install fastapi uvicorn
\end{lstlisting}

\section{Создание FastAPI-приложения}

Создайте файл \texttt{main.py} и добавьте:

\begin{lstlisting}[caption={Инициализация FastAPI-приложения}]
from fastapi import FastAPI

app = FastAPI(
    title="Agro Scoring API",
    description="API for agricultural risk scoring",
    version="1.0.0"
)
\end{lstlisting}

Здесь создаётся объект \texttt{app = FastAPI()}, который представляет наше API-приложение.

\section{Первый endpoint — /health}

Добавьте:

\begin{lstlisting}[caption={Endpoint /health}]
@app.get("/health")
def health():
    return {
        "status": "ok"
    }
\end{lstlisting}

Запустите приложение:

\begin{lstlisting}[language=bash]
uvicorn main:app --reload
\end{lstlisting}

Откройте \url{http://127.0.0.1:8000/health}. Ожидаемый результат:

\begin{lstlisting}[language=Python]
{"status": "ok"}
\end{lstlisting}

\section{Swagger UI}

Откройте \url{http://127.0.0.1:8000/docs}. FastAPI автоматически создаёт интерактивную документацию API. В ней должен появиться \texttt{GET /health}. Нажмите \texttt{Try it out}, затем \texttt{Execute}. Обратите внимание на Request URL, Response body, Response code и Response headers.

\section{OpenAPI}

FastAPI автоматически формирует описание API по стандарту OpenAPI. Откройте \url{http://127.0.0.1:8000/openapi.json} — вы увидите JSON-документ, описывающий ваш API. Упрощённо:

\begin{lstlisting}[caption={Фрагмент OpenAPI-спецификации}]
{
  "openapi": "3.x",
  "info": {
    "title": "Agro Scoring API"
  },
  "paths": {
    "/health": {}
  }
}
\end{lstlisting}

\section{Задание 1. Добавление информации о модели}

Создайте endpoint \texttt{GET /model-info}. Он должен возвращать:

\begin{lstlisting}[caption={Ответ /model-info}]
{
    "model_name": "agro-risk-model",
    "model_version": "1.0",
    "model_type": "risk-scoring",
    "status": "ready"
}
\end{lstlisting}

Пример реализации:

\begin{lstlisting}[caption={Endpoint /model-info}]
@app.get("/model-info")
def model_info():
    return {
        "model_name": "agro-risk-model",
        "model_version": "1.0",
        "model_type": "risk-scoring",
        "status": "ready"
    }
\end{lstlisting}

\textbf{Ответьте в отчёте:} чем \texttt{/health} отличается от \texttt{/model-info}?

\textbf{Ответ:} /health проверяет только работоспособность сервиса, а /model-info возвращает метаданные о ML-модели (имя, версию, тип, статус).

\section{Проектирование входных данных}

Для выполнения прогноза клиент должен передать сведения о хозяйстве. Например:

\begin{lstlisting}[caption={Пример входных данных}]
{
    "farm_id": "FARM-001",
    "region": "Krasnodar",
    "crop_type": "wheat",
    "area_ha": 2500,
    "temperature_avg": 24.3,
    "precipitation_mm": 320,
    "payment_delay_days": 15,
    "previous_defaults": 0,
    "debt": 1500000
}
\end{lstlisting}

FastAPI должен проверить эти данные до передачи модели.

\section{Pydantic}

FastAPI использует Pydantic для описания структуры данных. Добавьте:

\begin{lstlisting}[caption={Импорт Pydantic}]
from pydantic import BaseModel, Field
\end{lstlisting}

Создайте:

\begin{lstlisting}[caption={Модель FarmRequest}]
class FarmRequest(BaseModel):
    farm_id: str
    region: str
    crop_type: str
    area_ha: float = Field(gt=0)
    temperature_avg: float
    precipitation_mm: float = Field(ge=0)
    payment_delay_days: int = Field(ge=0)
    previous_defaults: int = Field(ge=0)
    debt: float = Field(ge=0)
\end{lstlisting}

Пояснения: \texttt{area\_ha: float = Field(gt=0)} означает тип данных \texttt{float}, значение строго больше 0.

\subsection*{Автоматическая валидация}

Если клиент передаст:

\begin{lstlisting}[caption={Некорректный запрос}]
{
    "area_ha": -200
}
\end{lstlisting}

FastAPI обнаружит ошибку ещё до запуска нашей функции инференса.

\section{Реализация модели}

В данной работе вместо модели машинного обучения будем использовать простую функцию, реализующую расчёт риска по некоторым финансовым и отраслевым показателям. Добавьте:

\begin{lstlisting}[caption={Функция расчёта риска}]
def calculate_risk(data: FarmRequest) -> float:
    score = 0.1
    if data.payment_delay_days > 30:
        score += 0.3
    if data.previous_defaults > 0:
        score += 0.3
    if data.debt > 5_000_000:
        score += 0.2
    if data.precipitation_mm < 100:
        score += 0.1
    return min(score, 1.0)
\end{lstlisting}

\section{Постпроцессинг}

Создайте:

\begin{lstlisting}[caption={Функция risk\_level}]
def risk_level(score: float) -> str:
    if score < 0.3:
        return "low"
    if score < 0.7:
        return "medium"
    return "high"
\end{lstlisting}

Добавьте рекомендации:

\begin{lstlisting}[caption={Функция recommendation}]
def recommendation(level: str) -> str:
    if level == "low":
        return "Standard review"
    if level == "medium":
        return "Additional check required"
    return "High risk. Manual review required"
\end{lstlisting}

\section{Задание 2. Реализуйте POST /predict}

Добавьте:

\begin{lstlisting}[caption={Импорт uuid}]
import uuid
\end{lstlisting}

Затем:

\begin{lstlisting}[caption={Endpoint /predict}]
@app.post("/predict")
def predict(request: FarmRequest):
    score = calculate_risk(request)
    level = risk_level(score)

    result = {
        "request_id": str(uuid.uuid4()),
        "farm_id": request.farm_id,
        "risk_score": score,
        "risk_level": level,
        "recommendation": recommendation(level),
        "model_version": "1.0"
    }
    return result
\end{lstlisting}

\section{Проверка через Swagger}

Откройте \texttt{/docs}. Выберите \texttt{POST /predict}. Введите:

\begin{lstlisting}[caption={Тестовый запрос}]
{
    "farm_id": "FARM-001",
    "region": "Krasnodar",
    "crop_type": "wheat",
    "area_ha": 2500,
    "temperature_avg": 24.3,
    "precipitation_mm": 320,
    "payment_delay_days": 45,
    "previous_defaults": 1,
    "debt": 6500000
}
\end{lstlisting}

Получите результат. Например:

\begin{lstlisting}[caption={Пример ответа}]
{
    "request_id": "7c8343c0-c912-4ab0-a56a-a4881894c7af",
    "farm_id": "FARM-001",
    "risk_score": 0.9,
    "risk_level": "high",
    "recommendation": "High risk. Manual review required",
    "model_version": "1.0"
}
\end{lstlisting}

\section{Модель ответа (Response Model)}

До настоящего момента мы описывали структуру входа, но не структуру ответа. Создайте:

\begin{lstlisting}[caption={Модель PredictionResponse}]
class PredictionResponse(BaseModel):
    request_id: str
    farm_id: str
    risk_score: float
    risk_level: str
    recommendation: str
    model_version: str
\end{lstlisting}

Измените endpoint:

\begin{lstlisting}[caption={Endpoint /predict с response\_model}]
@app.post(
    "/predict",
    response_model=PredictionResponse
)
def predict(request: FarmRequest):
    ...
\end{lstlisting}

Внешняя система должна получать ожидаемую структуру данных независимо от внутреннего устройства модели. Модель ответа позволяет контролировать формат ответа, автоматически документировать API, проверять данные ответа, скрывать ненужные внутренние поля и делать контракт между сервером и клиентом явным.

\section{Задание 3. Организовать хранение результатов}

Создадим временное хранилище в оперативной памяти (в виде словаря):

\begin{lstlisting}[caption={Хранилище прогнозов}]
predictions = {}
\end{lstlisting}

Измените \texttt{/predict}:

\begin{lstlisting}[caption={/predict с сохранением результата}]
@app.post(
    "/predict",
    response_model=PredictionResponse
)
def predict(request: FarmRequest):
    score = calculate_risk(request)
    level = risk_level(score)
    request_id = str(uuid.uuid4())

    result = {
        "request_id": request_id,
        "farm_id": request.farm_id,
        "risk_score": score,
        "risk_level": level,
        "recommendation": recommendation(level),
        "model_version": "1.0"
    }
    predictions[request_id] = result
    return result
\end{lstlisting}

\section{Path Parameters}

Предположим, необходимо получить конкретный прогноз по идентификатору. Запрос: \texttt{GET /predictions/abc-123}. Здесь \texttt{abc-123} является path parameter. FastAPI позволяет записать:

\begin{lstlisting}[caption={Path parameter}]
@app.get("/predictions/{request_id}")
def get_prediction(request_id: str):
    ...
\end{lstlisting}

Значение из URL попадёт в переменную \texttt{request\_id}.

\section{Задание 4. Получение прогноза}

Импортируйте:

\begin{lstlisting}[caption={Импорт HTTPException}]
from fastapi import HTTPException
\end{lstlisting}

Реализуйте:

\begin{lstlisting}[caption={Endpoint /predictions/\{request\_id\}}]
@app.get(
    "/predictions/{request_id}",
    response_model=PredictionResponse
)
def get_prediction(request_id: str):
    if request_id not in predictions:
        raise HTTPException(
            status_code=404,
            detail="Prediction not found"
        )
    return predictions[request_id]
\end{lstlisting}

\subsection*{HTTP коды ошибок}

Обратите внимание на 422. Попробуйте передать:

\begin{lstlisting}[caption={Некорректные данные}]
{
    "farm_id": "FARM-1",
    "region": "Krasnodar",
    "crop_type": "wheat",
    "area_ha": -100,
    "temperature_avg": 23,
    "precipitation_mm": 200,
    "payment_delay_days": 10,
    "previous_defaults": 0,
    "debt": 10000
}
\end{lstlisting}

FastAPI должен вернуть ошибку валидации. Обратите внимание на HTTP-код. Разберите структуру ответа.

\textbf{Ответьте:} почему такой запрос не должен доходить до ML-модели?

\textbf{Ответ:} Запрос с area\_ha = -100 не должен доходить до модели, потому что отрицательная площадь физически невозможна. Pydantic отсекает такие данные на этапе валидации (код 422), защищая модель от некорректного входа.

\section{Query Parameters}

Теперь реализуем получение списка прогнозов. Запрос: \texttt{GET /predictions?limit=10}. Здесь \texttt{limit=10} является query parameter. Пример:

\begin{lstlisting}[caption={Endpoint /predictions}]
@app.get("/predictions")
def get_predictions(limit: int = 10):
    values = list(predictions.values())
    return values[:limit]
\end{lstlisting}

\section{Задание 5. Фильтрация}

Добавьте возможность: \texttt{GET /predictions?risk\_level=high}. Например:

\begin{lstlisting}[caption={/predictions с фильтрацией}]
@app.get("/predictions")
def get_predictions(
    limit: int = 10,
    risk_level: str | None = None
):
    values = list(predictions.values())

    if risk_level is not None:
        values = [item for item in values if item["risk_level"] == risk_level]

    return values[:limit]
\end{lstlisting}

Протестируйте: \texttt{/predictions}, \texttt{/predictions?limit=2}, \texttt{/predictions?risk\_level=high}, \texttt{/predictions?risk\_level=medium\&limit=5}.

\section{Задание 6. Валидация Query Parameters}

Необходимо исключить запрос: \texttt{GET /predictions?limit=-100}. Импортируйте:

\begin{lstlisting}[caption={Импорт Query}]
from fastapi import Query
\end{lstlisting}

Измените параметр:

\begin{lstlisting}[caption={Валидация limit}]
limit: int = Query(
    default=10,
    ge=1,
    le=100
)
\end{lstlisting}

Теперь допустимый диапазон: $1 \le \text{limit} \le 100$.

\section{Задание 7. Проверка бизнес ограничений}

Pydantic проверяет формат данных, но некоторые ошибки невозможно определить только по типу значения, иногда необходима дополнительная проверка некоторых ограничений. Создайте:

\begin{lstlisting}[caption={Множество допустимых регионов}]
ALLOWED_REGIONS = {
    "Krasnodar",
    "Rostov",
    "Stavropol"
}
\end{lstlisting}

В \texttt{/predict} добавьте:

\begin{lstlisting}[caption={Бизнес-валидация региона}]
if request.region not in ALLOWED_REGIONS:
    raise HTTPException(
        status_code=400,
        detail="Unknown region"
    )
\end{lstlisting}

\section{Задание 8. Возвращаемые статусы}

Добавьте код ответа для \texttt{POST /predict}. Например:

\begin{lstlisting}[caption={Явное указание кода ответа}]
from fastapi import status

@app.post(
    "/predict",
    response_model=PredictionResponse,
    status_code=status.HTTP_201_CREATED
)
\end{lstlisting}

\textbf{Вопрос:} какой статус корректнее для инференса — «200 OK» или «201 Created»? Архитектор API должен выбрать семантику и использовать её последовательно.

\textbf{Ответ:} Если /predict создаёт сохраняемый ресурс (запись прогноза с request\_id), семантически корректнее 201 Created. В нашей реализации результат сохраняется в словаре predictions, поэтому выбран код 201.

\section{OpenAPI как контракт}

Откройте Swagger UI, теперь в документации автоматически отображаются \texttt{FarmRequest}, \texttt{PredictionResponse}, типы данных, обязательные поля, ограничения, HTTP endpoints и HTTP methods.

\section{Задание 9. Добавьте описание endpoints}

FastAPI позволяет сделать документацию более понятной. Например:

\begin{lstlisting}[caption={Описание /health}]
@app.get(
    "/health",
    summary="Health check",
    description="Returns current API status"
)
def health():
    ...
\end{lstlisting}

Для \texttt{/predict}:

\begin{lstlisting}[caption={Описание /predict}]
@app.post(
    "/predict",
    response_model=PredictionResponse,
    summary="Farm risk assessment",
    description="Runs inference of the agro scoring model"
)
def predict(request: FarmRequest):
    ...
\end{lstlisting}

Проверьте изменение Swagger UI.

\section{Задание 10. Тестирование API}

Последовательно выполните запросы и заполните таблицу тест-кейсов.

\begin{table}[H]
\centering
\caption{Таблица тест-кейсов}
\label{tab:test_cases}
\begin{tabular}{|c|p{5cm}|c|c|c|}
\hline
\textbf{№} & \textbf{Запрос} & \textbf{Ожидаемый HTTP-код} & \textbf{Полученный код} & \textbf{Результат} \\
\hline
1  & GET /health & 200 & 200 & успех \\
2  & GET /model-info & 200 & 200 & успех \\
3  & корректный POST /predict & 201 & 201 & успех \\
4  & area\_ha = -100 & 422 & 422 & успех \\
5  & неизвестный регион & 400 & 400 & успех \\
6  & существующий request\_id & 200 & 200 & успех \\
7  & неизвестный request\_id & 404 & 404 & успех \\
8  & limit = 2 & 200 & 200 & успех \\
9  & risk\_level = high & 200 & 200 & успех \\
10 & limit = -5 & 422 & 422 & успех \\
\hline
\end{tabular}
\end{table}

\section{Задание 11. Проанализируйте API как контракт}

Для каждого endpoint заполните таблицу.

\begin{table}[H]
\centering
\caption{Анализ API-контракта}
\label{tab:api_contract}
\begin{tabular}{|p{3cm}|p{1.5cm}|p{2.5cm}|p{3cm}|p{1.5cm}|p{2cm}|}
\hline
\textbf{Endpoint} & \textbf{Method} & \textbf{Input} & \textbf{Output} & \textbf{Success} & \textbf{Возможные ошибки} \\
\hline
/health & GET & --- & status & 200 & 500 \\
/model-info & GET & --- & model info & 200 & 500 \\
/predict & POST & FarmRequest & PredictionResponse & 201 & 400, 422, 500 \\
/predictions/\{id\} & GET & path id & PredictionResponse & 200 & 404 \\
/predictions & GET & query params & list & 200 & 422 \\
\hline
\end{tabular}
\end{table}

\section{Задание 12. Разделение приложения на модули}

Разделите приложение на несколько файлов:

\begin{lstlisting}[language=bash, caption={Структура проекта после рефакторинга}]
agro_api/
    main.py
    schemas.py
    model.py
    services.py
    storage.py
\end{lstlisting}

\textbf{schemas.py} содержит: \texttt{FarmRequest}, \texttt{PredictionResponse}.

\textbf{model.py} содержит: \texttt{calculate\_risk()}.

\textbf{services.py} содержит: \texttt{risk\_level()}, \texttt{recommendation()}, \texttt{is\_valid\_region()}, \texttt{build\_prediction()}.

\textbf{storage.py} содержит: \texttt{predictions}, \texttt{save\_prediction()}, \texttt{get\_prediction()}, \texttt{list\_predictions()}.

\textbf{main.py} содержит: инициализацию FastAPI и маршрутизацию endpoints.

Отчёт должен содержать:

\begin{enumerate}
    \item Цель работы.
    \item Краткое описание архитектуры API.
    \item Исходный код приложения: \url{https://github.com/dashakomarova2002-commits/HW_1_SISS}
    \item Реализованные endpoints: GET /health, GET /model-info, POST /predict, GET /predictions/\{request\_id\}, GET /predictions.
    \item Скриншот Swagger UI.
    \item Скриншот успешного /predict.
    \item Скриншот ошибки Pydantic.
    \item Скриншот 404 Not Found.
    \item Таблицу тест-кейсов.
    \item Таблицу API-контракта.
    \item В качестве вывода опишите разработанную архитектуру, желательно нарисовать компонентную схему архитектуры. Также сделайте краткий вывод о роли REST API в AI-системе.
\end{enumerate}

\section{Критерии оценки}

\begin{table}[H]
\centering
\caption{Критерии оценки}
\label{tab:criteria}
\begin{tabular}{|p{10cm}|c|}
\hline
\textbf{Критерий} & \textbf{Баллы} \\
\hline
Приложение запускается на FastAPI & 5 \\
Реализован /health & 5 \\
Реализован /model-info & 5 \\
Реализован /predict & 15 \\
Используется Pydantic & 10 \\
Реализована структурная валидация & 10 \\
Реализована бизнес-валидация & 10 \\
Реализован поиск по request\_id & 10 \\
Используются path/query parameters & 10 \\
Обрабатывается HTTP 404 & 5 \\
Корректно используются HTTP-коды & 5 \\
Проведено тестирование & 5 \\
Проанализирован OpenAPI-контракт & 5 \\
\hline
\textbf{Итого} & \textbf{100} \\
\hline
\end{tabular}
\end{table}

\section{Скриншоты работы приложения}

\begin{figure}[H]
    \centering
    \includegraphics[width=\textwidth]{01_swagger.png}
    \caption{Swagger UI со всеми endpoints}
\end{figure}

\begin{figure}[H]
    \centering
    \includegraphics[width=\textwidth]{02_predict_success.png}
    \caption{Успешный POST /predict (код 201)}
\end{figure}

\begin{figure}[H]
    \centering
    \includegraphics[width=\textwidth]{03_get_by_id.png}
    \caption{Получение прогноза по request\_id}
\end{figure}

\begin{figure}[H]
    \centering
    \includegraphics[width=\textwidth]{04_list.png}
    \caption{Список прогнозов с фильтрацией}
\end{figure}

\begin{figure}[H]
    \centering
    \includegraphics[width=\textwidth]{05_404.png}
    \caption{Ошибка 404 Not Found}
\end{figure}

\begin{figure}[H]
    \centering
    \includegraphics[width=\textwidth]{06_422.png}
    \caption{Ошибка валидации Pydantic (422)}
\end{figure}

\begin{figure}[H]
    \centering
    \includegraphics[width=\textwidth]{07_400.png}
    \caption{Ошибка бизнес-валидации региона (400)}
\end{figure}

\section*{Приложение А: Полный код приложения}

\begin{lstlisting}[caption={schemas.py}]
from pydantic import BaseModel, Field


class FarmRequest(BaseModel):
    farm_id: str
    region: str
    crop_type: str
    area_ha: float = Field(gt=0)
    temperature_avg: float
    precipitation_mm: float = Field(ge=0)
    payment_delay_days: int = Field(ge=0)
    previous_defaults: int = Field(ge=0)
    debt: float = Field(ge=0)


class PredictionResponse(BaseModel):
    request_id: str
    farm_id: str
    risk_score: float
    risk_level: str
    recommendation: str
    model_version: str
\end{lstlisting}

\begin{lstlisting}[caption={model.py}]
from schemas import FarmRequest


def calculate_risk(data: FarmRequest) -> float:
    score = 0.1
    if data.payment_delay_days > 30:
        score += 0.3
    if data.previous_defaults > 0:
        score += 0.3
    if data.debt > 5_000_000:
        score += 0.2
    if data.precipitation_mm < 100:
        score += 0.1
    return min(score, 1.0)
\end{lstlisting}

\begin{lstlisting}[caption={services.py}]
from schemas import FarmRequest
from model import calculate_risk

ALLOWED_REGIONS = {"Krasnodar", "Rostov", "Stavropol"}


def risk_level(score: float) -> str:
    if score < 0.3:
        return "low"
    if score < 0.7:
        return "medium"
    return "high"


def recommendation(level: str) -> str:
    if level == "low":
        return "Standard review"
    if level == "medium":
        return "Additional check required"
    return "High risk. Manual review required"


def is_valid_region(region: str) -> bool:
    return region in ALLOWED_REGIONS


def build_prediction(request: FarmRequest, request_id: str) -> dict:
    score = calculate_risk(request)
    level = risk_level(score)
    return {
        "request_id": request_id,
        "farm_id": request.farm_id,
        "risk_score": score,
        "risk_level": level,
        "recommendation": recommendation(level),
        "model_version": "1.0"
    }
\end{lstlisting}

\begin{lstlisting}[caption={storage.py}]
predictions: dict = {}


def save_prediction(request_id: str, result: dict) -> None:
    predictions[request_id] = result


def get_prediction(request_id: str) -> dict | None:
    return predictions.get(request_id)


def list_predictions(limit: int = 10, risk_level: str | None = None) -> list:
    values = list(predictions.values())
    if risk_level is not None:
        values = [item for item in values if item["risk_level"] == risk_level]
    return values[:limit]
\end{lstlisting}

\begin{lstlisting}[caption={main.py}]
import uuid

from fastapi import FastAPI, HTTPException, Query, status

from schemas import FarmRequest, PredictionResponse
from services import build_prediction, is_valid_region
from storage import save_prediction, get_prediction, list_predictions

app = FastAPI(
    title="Agro Scoring API",
    description="API for agricultural risk scoring",
    version="1.0.0"
)


@app.get("/health", summary="Health check")
def health():
    return {"status": "ok"}


@app.get("/model-info")
def model_info():
    return {
        "model_name": "agro-risk-model",
        "model_version": "1.0",
        "model_type": "risk-scoring",
        "status": "ready"
    }


@app.post(
    "/predict",
    response_model=PredictionResponse,
    status_code=status.HTTP_201_CREATED,
    summary="Farm risk assessment"
)
def predict(request: FarmRequest):
    if not is_valid_region(request.region):
        raise HTTPException(status_code=400, detail="Unknown region")

    request_id = str(uuid.uuid4())
    result = build_prediction(request, request_id)
    save_prediction(request_id, result)
    return result


@app.get(
    "/predictions/{request_id}",
    response_model=PredictionResponse
)
def get_prediction_endpoint(request_id: str):
    result = get_prediction(request_id)
    if result is None:
        raise HTTPException(status_code=404, detail="Prediction not found")
    return result


@app.get("/predictions")
def get_predictions(
    limit: int = Query(default=10, ge=1, le=100),
    risk_level: str | None = None
):
    return list_predictions(limit=limit, risk_level=risk_level)
\end{lstlisting}

\section*{Вывод}

В ходе работы спроектирован и реализован AI-сервис агроскоринга на FastAPI. Архитектура включает: API-слой (FastAPI), слой валидации (Pydantic-схемы FarmRequest и PredictionResponse), контур инференса (функция calculate\_risk и постпроцессинг risk\_level и recommendation) и хранилище результатов (словарь predictions в оперативной памяти).

REST API в AI-системе выполняет три ключевые роли: точка входа для всех внешних потребителей (CRM, портал, мобильное приложение), изолятор деталей реализации модели и контракт, фиксирующий форматы данных, коды ответов и правила обработки ошибок. Благодаря этому модель можно заменить или переобучить без изменения кода клиентов.

Компонентная схема:

\begin{center}
Клиент (CRM / веб-портал / мобильное приложение) \\
$\downarrow$ HTTP/JSON \\
FastAPI (main.py) \\
$\downarrow$ \\
Pydantic-схемы (schemas.py) $\rightarrow$ валидация \\
$\downarrow$ \\
Модель (model.py) $\rightarrow$ расчёт риска и постпроцессинг \\
$\downarrow$ \\
Хранилище (predictions) $\rightarrow$ ответ клиенту
\end{center}

\end{document}